% TerraPulse AI — Documentación Técnica y Científica: Simulación de Terremotos y Propagación de Ondas
% Compilar con: pdflatex ecuaciones_sismo.tex
\documentclass[11pt,a4paper]{article}
% Preámbulo común de la documentación LaTeX de TerraPulse AI.
% Se incluye con % Preámbulo común de la documentación LaTeX de TerraPulse AI.
% Se incluye con % Preámbulo común de la documentación LaTeX de TerraPulse AI.
% Se incluye con \input{preambulo} desde cada documento. Compilar con pdflatex (o latexmk -pdf).
\usepackage[utf8]{inputenc}
\usepackage[T1]{fontenc}
\usepackage{lmodern}
\usepackage[spanish,es-noshorthands,es-tabla,es-nodecimaldot]{babel}
\usepackage{amsmath,amssymb}
\usepackage[margin=2.3cm]{geometry}
\usepackage{array,booktabs,longtable}
\usepackage[normalem]{ulem}
\usepackage{xcolor}
\usepackage{enumitem}
\usepackage{listings}
\usepackage[hidelinks]{hyperref}

\setlength{\parindent}{0pt}
\setlength{\parskip}{0.5em}
\setlist{itemsep=2pt,topsep=2pt}
\renewcommand{\arraystretch}{1.25}
\Urlmuskip=0mu plus 1mu

% Símbolos frecuentes
\newcommand{\sz}{\ensuremath{\sigma^0}}
\newcommand{\gz}{\ensuremath{\gamma^0}}
\newcommand{\bz}{\ensuremath{\beta^0}}
\newcommand{\dsz}{\ensuremath{\Delta\sigma^0}}
\newcommand{\ok}{\ensuremath{\checkmark}}
\newcommand{\flecha}{\ensuremath{\rightarrow}}
\newcommand{\aprox}{\ensuremath{\approx}}
\newcommand{\por}{\ensuremath{\times}}

% Recuadro de nota (equivale a los ">" del Markdown)
\newenvironment{nota}
  {\begin{quote}\small\itshape}
  {\end{quote}}

% Bloques de código y comandos
\lstset{
  basicstyle=\ttfamily\small,
  breaklines=true,
  columns=fullflexible,
  keepspaces=true,
  frame=single,
  rulecolor=\color{black!30},
  backgroundcolor=\color{black!3},
  xleftmargin=4pt, xrightmargin=4pt,
  extendedchars=true,
  literate={á}{{\'a}}1 {é}{{\'e}}1 {í}{{\'i}}1 {ó}{{\'o}}1 {ú}{{\'u}}1
           {Á}{{\'A}}1 {É}{{\'E}}1 {Í}{{\'I}}1 {Ó}{{\'O}}1 {Ú}{{\'U}}1
           {ñ}{{\~n}}1 {Ñ}{{\~N}}1
}
 desde cada documento. Compilar con pdflatex (o latexmk -pdf).
\usepackage[utf8]{inputenc}
\usepackage[T1]{fontenc}
\usepackage{lmodern}
\usepackage[spanish,es-noshorthands,es-tabla,es-nodecimaldot]{babel}
\usepackage{amsmath,amssymb}
\usepackage[margin=2.3cm]{geometry}
\usepackage{array,booktabs,longtable}
\usepackage[normalem]{ulem}
\usepackage{xcolor}
\usepackage{enumitem}
\usepackage{listings}
\usepackage[hidelinks]{hyperref}

\setlength{\parindent}{0pt}
\setlength{\parskip}{0.5em}
\setlist{itemsep=2pt,topsep=2pt}
\renewcommand{\arraystretch}{1.25}
\Urlmuskip=0mu plus 1mu

% Símbolos frecuentes
\newcommand{\sz}{\ensuremath{\sigma^0}}
\newcommand{\gz}{\ensuremath{\gamma^0}}
\newcommand{\bz}{\ensuremath{\beta^0}}
\newcommand{\dsz}{\ensuremath{\Delta\sigma^0}}
\newcommand{\ok}{\ensuremath{\checkmark}}
\newcommand{\flecha}{\ensuremath{\rightarrow}}
\newcommand{\aprox}{\ensuremath{\approx}}
\newcommand{\por}{\ensuremath{\times}}

% Recuadro de nota (equivale a los ">" del Markdown)
\newenvironment{nota}
  {\begin{quote}\small\itshape}
  {\end{quote}}

% Bloques de código y comandos
\lstset{
  basicstyle=\ttfamily\small,
  breaklines=true,
  columns=fullflexible,
  keepspaces=true,
  frame=single,
  rulecolor=\color{black!30},
  backgroundcolor=\color{black!3},
  xleftmargin=4pt, xrightmargin=4pt,
  extendedchars=true,
  literate={á}{{\'a}}1 {é}{{\'e}}1 {í}{{\'i}}1 {ó}{{\'o}}1 {ú}{{\'u}}1
           {Á}{{\'A}}1 {É}{{\'E}}1 {Í}{{\'I}}1 {Ó}{{\'O}}1 {Ú}{{\'U}}1
           {ñ}{{\~n}}1 {Ñ}{{\~N}}1
}
 desde cada documento. Compilar con pdflatex (o latexmk -pdf).
\usepackage[utf8]{inputenc}
\usepackage[T1]{fontenc}
\usepackage{lmodern}
\usepackage[spanish,es-noshorthands,es-tabla,es-nodecimaldot]{babel}
\usepackage{amsmath,amssymb}
\usepackage[margin=2.3cm]{geometry}
\usepackage{array,booktabs,longtable}
\usepackage[normalem]{ulem}
\usepackage{xcolor}
\usepackage{enumitem}
\usepackage{listings}
\usepackage[hidelinks]{hyperref}

\setlength{\parindent}{0pt}
\setlength{\parskip}{0.5em}
\setlist{itemsep=2pt,topsep=2pt}
\renewcommand{\arraystretch}{1.25}
\Urlmuskip=0mu plus 1mu

% Símbolos frecuentes
\newcommand{\sz}{\ensuremath{\sigma^0}}
\newcommand{\gz}{\ensuremath{\gamma^0}}
\newcommand{\bz}{\ensuremath{\beta^0}}
\newcommand{\dsz}{\ensuremath{\Delta\sigma^0}}
\newcommand{\ok}{\ensuremath{\checkmark}}
\newcommand{\flecha}{\ensuremath{\rightarrow}}
\newcommand{\aprox}{\ensuremath{\approx}}
\newcommand{\por}{\ensuremath{\times}}

% Recuadro de nota (equivale a los ">" del Markdown)
\newenvironment{nota}
  {\begin{quote}\small\itshape}
  {\end{quote}}

% Bloques de código y comandos
\lstset{
  basicstyle=\ttfamily\small,
  breaklines=true,
  columns=fullflexible,
  keepspaces=true,
  frame=single,
  rulecolor=\color{black!30},
  backgroundcolor=\color{black!3},
  xleftmargin=4pt, xrightmargin=4pt,
  extendedchars=true,
  literate={á}{{\'a}}1 {é}{{\'e}}1 {í}{{\'i}}1 {ó}{{\'o}}1 {ú}{{\'u}}1
           {Á}{{\'A}}1 {É}{{\'E}}1 {Í}{{\'I}}1 {Ó}{{\'O}}1 {Ú}{{\'U}}1
           {ñ}{{\~n}}1 {Ñ}{{\~N}}1
}

\setlength{\parskip}{0.4em}

\title{\textbf{TerraPulse AI}\\[6pt]\Large Manual Técnico y Físico-Matemático: Simulación de Terremotos, Propagación de Ondas Elásticas, Atenuación GMPE y Zonificación de Riesgo Sísmico}
\author{Ingeniería de Sistemas y Geociencias Computacionales -- TerraPulse AI}
\date{17 de noviembre de 2026}

\begin{document}
\maketitle

\begin{abstract}
Este documento constituye la memoria técnica integral de la arquitectura, ecuaciones y funcionamiento del \textbf{Simulador de Terremotos de TerraPulse AI}. Se detalla la formulación física basada en la elastodinámica continua (ecuación de Navier), la determinación de tiempos de viaje para ondas primarias ($P$) y secundarias ($S$), la estimación de la Aceleración Pico del Terreno ($PGA$) mediante Leyes de Predicción de Movimiento del Suelo (\emph{Ground Motion Prediction Equations}, GMPE) calibradas para la zona de subducción e intraplaca de México (Ordaz et al., García et al., Boore-Atkinson), la amplificación estratigráfica por velocidad de onda de corte en 30 metros ($V_{s30}$), la correlación con la escala Mercalli Modificada ($MMI$) de daño estructural, la generación analítica del sismograma sintético $a(t, D)$ con efecto de resonancia de cuenca (Valle de México) y el algoritmo numérico de bisección para la zonificación de 4 anillos concéntricos de isosisma georreferenciados en MapLibre GL.
\end{abstract}

\tableofcontents
\newpage

\section{Introducción y Visión General del Sistema}

TerraPulse AI es una plataforma de análisis geoespacial y computación de riesgos naturales. Su módulo de simulación de terremotos permite a analistas de protección civil, investigadores y tomadores de decisiones simular en tiempo real escenarios de ruptura sísmica, evaluar el impacto destructivo del movimiento del suelo y estimar con precisión el tiempo disponible de alerta temprana.

El simulador opera combinando tres pilares científicos:
\begin{enumerate}
  \item \textbf{Física del medio continuo y elastodinámica:} Cálculo riguroso de tiempos de arribo de frentes de onda longitudinales ($P$) y transversales ($S$) a partir de las propiedades elásticas de la corteza continental.
  \item \textbf{Modelos empíricos y leyes de atenuación GMPE:} Estimación probabilística de aceleraciones espectrales calibradas con registros históricos de la Trinchera Mesoamericana.
  \item \textbf{Cómputo geoespacial geodésico y series temporales:} Inversión numérica de radios de afectación y generación de trazas sintéticas de acelerogramas $a(t)$ con respuesta de sitio.
\end{enumerate}

\section{Arquitectura y Flujo de Procesamiento del Simulador}

El ciclo de cómputo y renderizado del simulador opera de forma reactiva e interactiva siguiendo las etapas descritas a continuación:

\begin{enumerate}
  \item \textbf{Captura de Parámetros de Ruptura (Hipocentro y Mecanismo):}
    El usuario define las coordenadas del epicentro $(\lambda_e, \phi_e)$ mediante un clic directo en el mapa o seleccionando un sismo histórico (p.~ej., Michoacán 1985 $M_w 8.1$, Puebla-Morelos 2017 $M_w 7.1$ o Tehuantepec 2017 $M_w 8.2$). Se configuran la magnitud de momento $M_w$, la profundidad focal $h$ (km), el mecanismo de falla (subducción interplaca, falla normal intraplaca o cortical superficial) y la litología superior ($V_{s30}$).

  \item \textbf{Cálculo Geodésico y Distancia Hipocentral:}
    Para cualquier punto del dominio, se evalúa la distancia ortodrómica sobre el esferoide terrestre y la distancia hipocentral $R = \sqrt{D^2 + h^2}$.

  \item \textbf{Evaluación del Campo de Movimiento del Suelo ($PGA$ y $MMI$):}
    Se evalúa la ecuación GMPE considerando atenuación geométrica $\sim 1/R$, amortiguamiento inelástico $e^{-\gamma R}$ y el factor de amplificación de sitio $S_{\text{sitio}}$. El valor pico resultante se transforma a la escala Mercalli Modificada ($MMI$) para cuantificar daños en infraestructura.

  \item \textbf{Inversión Numérica de Radios de Isosisma (Bisección):}
    Mediante un algoritmo de bisección se resuelven las 4 distancias epicentrales límites para los umbrales de aceleración destructiva ($\ge 0.34\,g$), daño severo ($0.12\,g$), percepción moderada ($0.03\,g$) y percepción leve ($0.005\,g$).

  \item \textbf{Generación de Polígonos Geodésicos Vectoriales:}
    A partir de cada radio $R_{\text{zona}}$, se computa una trayectoria circular geodésica de 64 nodos sobre el elipsoide WGS84, ensamblando capas GeoJSON proyectadas dinámicamente en MapLibre GL.

  \item \textbf{Síntesis Temporal del Sismograma $a(t, D)$:}
    Se genera la serie continua de aceleración en función del tiempo para una estación receptora elegida por el usuario, mostrando la onda $P$, la ventana de advertencia temprana $\Delta t = t_S - t_P$, el impacto cortante de la onda $S$ y la coda con resonancia de estrato blando.
\end{enumerate}

\section{Ecuaciones de Fuente y Geometría Focal}

El momento sísmico escalar $M_0$ (en dyn$\cdot$cm) representa la energía liberada en la falla:
\begin{equation}
  M_w = \frac{2}{3} \log_{10}(M_0) - 10.7 \iff M_0 = 10^{\frac{3}{2}(M_w + 10.7)}.
  \label{eq:kanamori}
\end{equation}

Para una estación en $(\lambda, \phi)$ y epicentro en $(\lambda_e, \phi_e)$, la distancia epicentral superficial $D$ se obtiene por la ley de semiversenos (Haversine):
\begin{equation}
  D = 2 R_T \arcsin \left( \sqrt{\sin^2\left(\frac{\Delta \phi}{2}\right) + \cos(\phi_e)\cos(\phi)\sin^2\left(\frac{\Delta \lambda}{2}\right)} \right),
  \label{eq:distancia_epicentral}
\end{equation}
donde $R_T \approx 6371\,\text{km}$, $\Delta \phi = \phi - \phi_e$ y $\Delta \lambda = \lambda - \lambda_e$.

La distancia hipocentral tridimensional $R$ desde la ruptura hasta la superficie es:
\begin{equation}
  R = \sqrt{D^2 + h^2}.
  \label{eq:distancia_hipocentral}
\end{equation}

\section{Elastodinámica Continua y Tiempos de Viaje de Ondas}

La propagación de deformaciones mecánicas en la corteza terrestre está regida por la ecuación elástica de Navier:
\begin{equation}
  \rho \frac{\partial^2 \mathbf{u}}{\partial t^2} = (\lambda + \mu)\nabla(\nabla \cdot \mathbf{u}) + \mu \nabla^2 \mathbf{u},
  \label{eq:navier}
\end{equation}
donde $\mathbf{u}(\mathbf{x}, t)$ es el vector desplazamiento, $\rho$ es la densidad de la roca, y $\lambda, \mu$ son los coeficientes elásticos de Lamé ($\mu$ módulo de corte o rigidez).

Aplicando el teorema de descomposición de Helmholtz $\mathbf{u} = \nabla \Phi + \nabla \times \mathbf{\Psi}$:
\begin{enumerate}
  \item \textbf{Potencial escalar $\Phi$ (Onda P, compresional / longitudinal):}
    \begin{equation}
      \nabla^2 \Phi = \frac{1}{v_p^2}\frac{\partial^2 \Phi}{\partial t^2}, \qquad v_p = \sqrt{\frac{\lambda + 2\mu}{\rho}} \approx 6.0\,\text{km/s}.
    \end{equation}
  \item \textbf{Potencial vectorial $\mathbf{\Psi}$ (Onda S, cizalla / transversal):}
    \begin{equation}
      \nabla^2 \mathbf{\Psi} = \frac{1}{v_s^2}\frac{\partial^2 \mathbf{\Psi}}{\partial t^2}, \qquad v_s = \sqrt{\frac{\mu}{\rho}} \approx 3.5\,\text{km/s}.
    \end{equation}
\end{enumerate}

Para un sólido de Poisson ($\nu = 0.25 \implies \lambda = \mu$), la razón cinemática satisface:
\begin{equation}
  \frac{v_p}{v_s} = \sqrt{\frac{\lambda + 2\mu}{\mu}} = \sqrt{3} \approx 1.732.
\end{equation}

\subsection{Tiempos de Llegada y Ventana Operativa de Alerta Temprana}

Los tiempos de viaje desde el foco hasta el sitio receptor son:
\begin{equation}
  t_P(R) = \frac{R}{v_p} = \frac{\sqrt{D^2 + h^2}}{6.0}, \qquad t_S(R) = \frac{R}{v_s} = \frac{\sqrt{D^2 + h^2}}{3.5}.
\end{equation}

El margen temporal entre el arribo de la onda $P$ (no destructiva) y el tren cortante destructivo $S$ define la ventana de alerta temprana:
\begin{equation}
  \Delta t(D) = t_S - t_P = \sqrt{D^2 + h^2} \left( \frac{1}{3.5} - \frac{1}{6.0} \right) \approx 0.119 \cdot \sqrt{D^2 + h^2} \quad \text{(segundos)}.
\end{equation}

Dado que los sensores sismológicos costeros (SASMEX) transmiten por enlace de radio a la velocidad de la luz $c \approx 3\times 10^5\,\text{km/s} \gg v_s$, el tiempo de oportunidad útil para alertar a una metrópoli ubicada a $D \approx 320\,\text{km}$ (como la Ciudad de México respecto a Guerrero o Michoacán) es:
\begin{equation}
  T_{\text{alerta}} \approx \frac{\sqrt{320^2 + 15^2}}{3.5} - t_{\text{proceso}} \approx 91.5\,\text{s} - 3.0\,\text{s} \approx 88.5\,\text{s} \quad (\text{tiempo total desde el origen}).
\end{equation}
La población percibe la alerta entre 40 y 60 segundos antes del movimiento fuerte.

\section{Leyes de Atenuación del Movimiento del Suelo (GMPE)}

El parámetro físico principal para ingeniería estructural es la Aceleración Pico del Terreno (\emph{Peak Ground Acceleration}, $PGA$). El simulador implementa el modelo regional adaptado a México (Ordaz, García, Boore-Atkinson):
\begin{equation}
  \ln(PGA) = c_1 + c_2 M_w - c_3 \ln(R) - c_4 R + S_{\text{sitio}},
  \label{eq:gmpe}
\end{equation}
donde $PGA$ está expresado en fracciones de la gravedad terrestre ($g$), con $1\,g = 980.665\,\text{cm/s}^2$ (Gals).

\subsection{Coeficientes según Mecanismo Sismotectónico}

\begin{itemize}
  \item \textbf{Subducción Interplaca (Trinchera Mesoamericana):}
    $c_1 = -1.85$, $c_2 = 0.92$, $c_3 = 1.15$, $c_4 = 0.0032\,\text{km}^{-1}$.
  \item \textbf{Intraplaca en Placa de Cocos Subducida (Sismos Profundos):}
    $c_1 = -1.55$, $c_2 = 0.96$, $c_3 = 1.25$, $c_4 = 0.0040\,\text{km}^{-1}$ (mayor contenido de alta frecuencia y menor atenuación por profundidad focal de 40 a 70 km).
  \item \textbf{Cortical Somero (Fallas Locales):}
    $c_1 = -1.95$, $c_2 = 0.88$, $c_3 = 1.05$, $c_4 = 0.0028\,\text{km}^{-1}$.
\end{itemize}

\subsection{Efecto de Sitio y Clasificación $V_{s30}$}

La amplificación estratigráfica $S_{\text{sitio}}$ depende de la velocidad promedio de onda de corte en los primeros 30 metros de estrato:
\begin{equation}
  S_{\text{sitio}} = b_{\text{site}} \ln\left( \frac{V_{s30}}{V_{\text{ref}}} \right), \quad V_{\text{ref}} = 760\,\text{m/s}, \quad b_{\text{site}} = -0.45.
\end{equation}

\begin{itemize}
  \item \textbf{Roca Firme ($V_{s30} \ge 760\,\text{m/s}$):} $S_{\text{sitio}} \approx 0.0$.
  \item \textbf{Suelo Firme / Depósitos Aluviales ($V_{s30} \approx 450\,\text{m/s}$):} Amplificación moderada ($\sim 1.3\times$).
  \item \textbf{Suelo Blando Lacustre ($V_{s30} \approx 100\,\text{m/s}$, Zona de Lago de la CDMX):} $S_{\text{sitio}} \approx 0.91$, lo que cuadruplica el valor de aceleración espectral debido al fenómeno de resonancia por bajo contraste de impedancias.
\end{itemize}

\section{Escala de Intensidad Mercalli Modificada (\texorpdfstring{$MMI$}{MMI})}

Para traducir la aceleración cuantitativa a impacto perceptual y daños estructurales, se aplica la correlación de Worden et al. y Wald et al.:
\begin{equation}
  MMI(PGA_{\text{cgs}}) =
  \begin{cases}
    2.20 \log_{10}(PGA_{\text{cgs}}) + 1.00, & \text{si } PGA_{\text{cgs}} \le 4.22\,\text{cm/s}^2 \\[6pt]
    3.66 \log_{10}(PGA_{\text{cgs}}) - 1.66, & \text{si } PGA_{\text{cgs}} > 4.22\,\text{cm/s}^2
  \end{cases}
  \label{eq:mmi}
\end{equation}
donde $PGA_{\text{cgs}} = PGA_g \times 980.665$.

\section{Zonificación de Riesgo Sísmico y Algoritmo de Bisección}

El simulador define 4 categorías operativas de riesgo sísmico:

\begin{table}[ht!]
\centering
\small
\begin{tabular}{@{}llccp{6.8cm}@{}}
\toprule
\textbf{Nivel de Riesgo} & \textbf{Color UI} & \textbf{Intensidad $MMI$} & \textbf{Umbral $PGA$} & \textbf{Efectos Estructurales y Poblacionales} \\
\midrule
\textbf{Zona Destructiva} & Rojo (\#ef4444) & $\ge \text{VIII}$ & $\ge 0.34\,g$ & Daño severo y colapso de edificios; rotura masiva de redes hidráulicas y gasoductos; pánico total. \\
\textbf{Zona de Daño Severo} & Naranja (\#f97316) & $\text{VI - VII}$ & $0.12 - 0.34\,g$ & Daño moderado en muros y cornisas; fisuras estructurales; objetos pesados desplazados. \\
\textbf{Zona Moderada} & Ámbar (\#f59e0b) & $\text{IV - V}$ & $0.03 - 0.12\,g$ & Sentido ampliamente por toda la población; balanceo de lámparas y vajillas; sin daño estructural. \\
\textbf{Zona de Percepción} & Azul (\#3b82f6) & $\text{II - III}$ & $0.005 - 0.03\,g$ & Sentido únicamente por personas en reposo o edificios elevados; límite instrumental. \\
\bottomrule
\end{tabular}
\caption{Criterios cuantitativos de zonificación de riesgo sísmico en TerraPulse AI.}
\label{tab:riesgo}
\end{table}

Para cada umbral $PGA_{\text{umbral}}$, el radio de afectación $R_{\text{zona}}$ se calcula numéricamente mediante el método de bisección en el intervalo $[0, 1500\,\text{km}]$ con 30 iteraciones (precisión $< 0.01\,\text{km}$):
\begin{equation}
  f(R) = PGA(R, M_w, h, \dots) - PGA_{\text{umbral}} = 0.
\end{equation}

A partir de $R_{\text{zona}}$, se genera un anillo geodésico sobre la esfera terrestre evaluando 64 vértices mediante trigonometría esférica de punto directo:
\begin{align}
  \phi_i &= \arcsin\left( \sin\phi_e \cos\delta + \cos\phi_e \sin\delta \cos\theta_i \right), \\
  \lambda_i &= \lambda_e + \text{atan2}\left( \sin\theta_i \sin\delta \cos\phi_e, \; \cos\delta - \sin\phi_e \sin\phi_i \right),
\end{align}
donde $\delta = R_{\text{zona}} / R_T$ y $\theta_i = \frac{2\pi i}{64}$.

\section{Formulación Analítica del Sismograma Sintético \texorpdfstring{$a(t, D)$}{a(t, D)}}

Para reproducir fielmente la señal oscilatoria que registra una estación sísmica a distancia epicentral $D$, el simulador genera la serie temporal:
\begin{equation}
  a(t, D) = a_{\text{ruido}}(t) + a_P(t, D) + a_S(t, D).
  \label{eq:sismograma}
\end{equation}

\subsection{Paquete de Onda Primaria (\texorpdfstring{$P$}{P})}
Arriba en $t_P = \sqrt{D^2 + h^2} / 6.0$. Posee una vibración compresional de alta frecuencia ($f_P = 4.5\,\text{Hz}$) y baja amplitud ($\sim 22\%$ de $PGA$):
\begin{equation}
  a_P(t, D) = 0.22 \cdot PGA(D) \cdot \sin(2\pi f_P (t - t_P)) \cdot \cos(0.3\pi f_P (t - t_P)) \cdot \exp\left(-\frac{t - t_P}{3.2}\right) \cdot \mathcal{H}(t - t_P).
\end{equation}

\subsection{Paquete de Onda Secundaria (\texorpdfstring{$S$}{S}) y Ondas Superficiales}
Arriba en $t_S = \sqrt{D^2 + h^2} / 3.5$. Concentra el 80\% de la energía cinética destructiva, alcanzando el valor pico $PGA(D)$:
\begin{equation}
  a_S(t, D) = PGA(D) \cdot E_S(t - t_S) \cdot \Big[ 0.82 \sin(2\pi f_S (t - t_S)) + 0.18 \sin(4\pi f_S (t - t_S) + 0.6) \Big] \cdot \mathcal{H}(t - t_S),
\end{equation}
con la envolvente $E_S(\tau)$:
\begin{equation}
  E_S(\tau) =
  \begin{cases}
    \dfrac{\tau}{1.8}, & 0 \le \tau \le 1.8\,\text{s} \\[8pt]
    \exp\left(-\dfrac{\tau - 1.8}{\tau_{\text{coda}}}\right), & \tau > 1.8\,\text{s}.
  \end{cases}
\end{equation}

\subsection{Efecto de Resonancia en Cuenca Lacustre (CDMX)}
En depósitos lacustres blandos ($V_{s30} < 180\,\text{m/s}$), el contraste de impedancia acústica entre el basamento rocoso y las arcillas de alta compresibilidad atrapa las ondas sísmicas. En el simulador, esto se modela reduciendo la frecuencia fundamental a $f_S = 0.52\,\text{Hz}$ (período característico $T_0 \approx 1.9 - 2.0\,\text{s}$) y elevando el tiempo de decaimiento a $\tau_{\text{coda}} = 28.0\,\text{s}$. Esto reproduce fielmente las sacudidas armónicas de gran persistencia observadas en 1985 y 2017.

\section{Presets Históricos Mexicanos Precalibrados}

El sistema incluye 6 escenarios históricos calibrados para validación analítica:
\begin{enumerate}
  \item \textbf{1985 -- Michoacán / Costa del Pacífico ($M_w 8.1$, $h=15\,\text{km}$):} Ruptura de subducción en la Trinchera Mesoamericana. Generó severa resonancia y colapso estructural en la Ciudad de México a 350 km de distancia.
  \item \textbf{2017 -- Puebla / Morelos ($M_w 7.1$, $h=57\,\text{km}$):} Falla normal intraplaca profunda en el slab de Cocos. Altas aceleraciones destructivas de alta frecuencia en Morelos, Puebla y el Valle de México.
  \item \textbf{2017 -- Golfo de Tehuantepec ($M_w 8.2$, $h=45\,\text{km}$):} El sismo de mayor magnitud registrado en México en un siglo. Falla normal intraplaca con extensa zona de ruptura en Oaxaca y Chiapas.
  \item \textbf{2021 -- Acapulco / Guerrero ($M_w 7.1$, $h=10\,\text{km}$):} Subducción somera en la brecha de Guerrero con fuerte aceleración en la costa y activación de alerta temprana en CDMX.
  \item \textbf{2022 -- Coalcomán / Michoacán ($M_w 7.7$, $h=15\,\text{km}$):} Ruptura interplaca con afectaciones en Colima, Michoacán y alerta sísmica en la zona centro.
  \item \textbf{2010 -- Mexicali / Baja California ($M_w 7.2$, $h=10\,\text{km}$):} Falla cortical de rumbo lateral (sistema Laguna Salada / San Andrés).
\end{enumerate}

\section{Mapeo Sinóptico entre Ecuaciones y Funciones del Código}

\begin{table}[ht!]
\centering
\small
\begin{tabular}{@{}lllp{5.5cm}@{}}
\toprule
\textbf{Magnitud Física} & \textbf{Símbolo / Ecuación} & \textbf{Unidad} & \textbf{Función en TypeScript (\texttt{seismic.ts})} \\
\midrule
Distancia Epicentral & $D = 2 R_T \arcsin(\dots)$ & km & Trigonometría esférica Haversine \\
Distancia Hipocentral & $R = \sqrt{D^2 + h^2}$ & km & \texttt{Math.sqrt(dist*dist + depthKm*depthKm)} \\
Aceleración Pico & $\ln(PGA) = c_1 + c_2 M_w - c_3 \ln R - c_4 R + S$ & $g$ & \texttt{computePga\_g(...)} \\
Aceleración en CGS & $PGA_{\text{cgs}} = PGA_g \times 980.665$ & $\text{cm/s}^2$ & \texttt{pga\_g * 980.665} \\
Escala Mercalli & $MMI = 3.66 \log_{10}(PGA_{\text{cgs}}) - 1.66$ & I a X+ & \texttt{pgaToMmi(pga\_g)} \\
Tiempo de Onda P & $t_P = R / 6.0$ & s & \texttt{tP = hypoDist / 6.0} \\
Tiempo de Onda S & $t_S = R / 3.5$ & s & \texttt{tS = hypoDist / 3.5} \\
Ventana Alerta SASMEX & $\Delta t = t_S - t_P$ & s & \texttt{Math.max(0, tS - tP)} \\
Sismograma $a(t)$ & $a(t) = a_{\text{ruido}} + a_P(t) + a_S(t)$ & $g$, Gal & \texttt{generateSyntheticSeismogram(...)} \\
Radios de Afectación & $\text{Bisección: } PGA(R_{\text{zona}}) = PGA_{\text{umbral}}$ & km & \texttt{solveDistanceForPga(...)} \\
Polígono Isosisma & Proyección geodésica $64$ nodos & LngLat[] & \texttt{generateGeodesicCircle(center, rKm, 64)} \\
\bottomrule
\end{tabular}
\caption{Correspondencia exacta entre el modelo físico analítico y las rutinas de cálculo del código fuente.}
\label{tab:mapeo}
\end{table}

\section{Conclusiones}

La integración de la elastodinámica continua con las leyes regionales de atenuación GMPE y la resonancia estratigráfica permite que el Simulador de Terremotos de TerraPulse AI proporcione una herramienta analítica confiable, científicamente rigurosa y visualmente accesible en tiempo real tanto en cartografía vectorial como en análisis de series temporales.

\end{document}
